\chapter{Lo que la dopamina dice en realidad}
% versión completa: chapters/07_dofaalt.tex

En el capítulo anterior la dopamina era sencilla: un solo número, “mejor o peor de lo esperado”. Es una primera aproximación y funciona bien. Pero los experimentos de los últimos años mostraron que por ese “cable” viaja más que un número. Este capítulo trata de dónde el cuadro se complica y dónde los científicos todavía discuten.

\section*{Optimistas y pesimistas}

Las neuronas de dopamina no son todas iguales. Unas reaccionan más a las sorpresas agradables que a las desagradables: son las optimistas. Otras, al revés: las pesimistas. Si cada neurona aprende con su propia regla, la optimista aprenderá cómo será la recompensa en el mejor caso, y la pesimista, en el peor. Juntas no guardan una sola expectativa media, sino toda la dispersión de resultados posibles, como un corredor de apuestas que no solo conoce al favorito, sino las probabilidades de cada uno. La idea concuerda con los experimentos; para qué quiere el cerebro toda la dispersión es, por ahora, una hipótesis.

\pencilfig{7}{\pencilart{7}}{Cada neurona de dopamina guarda su propio punto de la curva de recompensas posibles: juntas recuerdan toda la curva.}

\section*{No solo “mejor”, también “importante”}

Parte de las neuronas de dopamina se encienden también ante sucesos desagradables pero importantes: un ruido fuerte, un peligro. Parece que por la red de dopamina no solo viaja el signo de la sorpresa, sino también su relevancia: una señal de “presta atención”.

\section*{Dos señales en un cable}

Los estallidos rápidos enseñan. Pero el nivel medio lento de dopamina, que cambia en minutos, parece hablar de otra cosa: de lo rico que es ahora el entorno. Si hay mucha recompensa alrededor, el tiempo vale, y el animal actúa más rápido. Los experimentos confirman que el nivel de dopamina está ligado al vigor de las acciones. Un ingeniero lo llamaría separación de señales por frecuencia: las altas frecuencias llevan una cosa y las bajas, otra.

\section*{Una rampa hacia la meta}

Cuando el animal se acerca a la recompensa, la dopamina a menudo no estalla, sino que sube de forma suave. Una explicación: es la expectativa misma, una señal de “la meta está cerca, apura”. Otra conserva la idea de sorpresa: de lejos el animal evalúa la situación de forma borrosa, y al acercarse, cada vez con más precisión, y cada precisión nueva es una pequeña sorpresa agradable. Un experimento elegante en un pasillo virtual, donde se teletransportaba al animal más cerca de la meta, apoya la segunda: la dopamina respondió con un estallido, como corresponde a una sorpresa.

\section*{Mirar hacia atrás}

Hay también una alternativa audaz: la dopamina no responde a la pregunta “¿qué va a pasar?”, sino a “¿qué causó la recompensa?”; mira al pasado y no al futuro. Los experimentos en los que las dos teorías predicen cosas distintas se contradicen por ahora entre sí. Es un debate científico vivo.

\section*{Un vector en lugar de un número}

Por último, la dopamina reacciona también a un cambio de \emph{sabor} de la recompensa cuando su valor no cambió, y ayuda a aprender sin recompensa alguna. Parece que no es un solo error, sino todo un haz de errores: uno por cada rasgo de lo ocurrido.

La conclusión del capítulo: la mayoría de las teorías nuevas no anulan el cuadro sencillo, sino que lo amplían. La dopamina no es un solo cable, sino un pequeño mazo de varios. Solo “mirar hacia atrás” lo discute en serio.

\fullversion{7}
